% ============================================================
%  Method: Face Restoration via Image-to-Video Generation
%  using TurboDiffusion (Wan2.2-A14B)
%
%  Drop this into main.tex as a \section or \subsection
% ============================================================

\subsection{Face restoration via image-to-video generation}
\label{subsec:i2v_method}

Rather than restoring a degraded face image into a single enhanced
frame, we propose to exploit an \emph{image-to-video} (I2V) diffusion
model as the restoration backbone.  The core idea is that a video
generation model conditioned on the degraded input will produce a
temporal sequence of $T$ frames; among these frames, at least one is
likely to exhibit a high-fidelity, near-frontal face that better
preserves the subject's identity than any single-image method can
achieve.

%------------------------------------------------------------
\subsubsection{TurboDiffusion framework}
\label{subsubsec:turbodiffusion}

We employ TurboDiffusion~\cite{TurboDiffusion}, a video generation
acceleration framework built upon the Wan2.2-A14B
architecture---a 14-billion-parameter image-to-video diffusion
model.  TurboDiffusion achieves $100$--$200 \times$ end-to-end
speedup compared to the original Wan~2.2 model by combining
three synergistic acceleration techniques:

\begin{enumerate}
    \item \textbf{rCM (Rectified Consistency Model)
          distillation}~\cite{rCM}.\;
          The standard 50--100 denoising steps of the base model are
          compressed to only \textbf{4~steps} via score-regularized
          continuous-time consistency training.  Two complementary
          checkpoints are trained for different noise regimes:
          \begin{itemize}
              \item \emph{Low-noise model} ($\bm{\theta}_{\text{low}}$):
                    handles the low-noise range
                    (timestep boundary $\leq 0.9$).
              \item \emph{High-noise model} ($\bm{\theta}_{\text{high}}$):
                    handles the high-noise range
                    (timestep boundary $> 0.9$).
          \end{itemize}
          During inference, the scheduler switches from
          $\bm{\theta}_{\text{high}}$ to $\bm{\theta}_{\text{low}}$
          when the current timestep crosses the boundary.

    \item \textbf{SageSLA (Sparse-Linear Attention)}~\cite{SLA}.\;
          Full quadratic self-attention $\mathcal{O}(n^{2})$ in the
          DiT backbone is replaced by a sparse-linear variant that
          selects only the top-$k$ attention entries ($k=0.1$),
          reducing attention computation by approximately $10\times$
          while preserving generation quality.

    \item \textbf{INT8 weight quantization}.\;
          All linear projection layers in the DiT are quantized to
          INT8, halving memory bandwidth and enabling inference on
          consumer GPUs (e.g., NVIDIA RTX~5090 with ${<}24$\,GB VRAM).
\end{enumerate}

The resulting pipeline generates a 720p, 81-frame video in
approximately \textbf{38 seconds} on a single RTX~5090, compared
to $\sim$4549\,s for the unaccelerated Wan~2.2 model---a
$120\times$ speedup.

%------------------------------------------------------------
\subsubsection{Conditional video generation}
\label{subsubsec:conditional_gen}

Given a degraded input image $B_i(\sigma)$ and a text prompt
$\mathcal{P}$, the I2V model generates a video clip
$\mathbf{V} = \{I_t\}_{t=1}^{T}$ with $T = 81$ frames.
The generation follows an ODE (Ordinary Differential Equation)
sampling trajectory, which yields deterministic and sharper outputs
compared to stochastic SDE sampling.

We design three severity-adaptive text prompts, each tailored to
guide the generation process according to the degree of input
degradation:

\paragraph{Mild blur ($\sigma \in \{3, 5\}$).}
\begin{quote}\small\itshape
``Ultra-high definition macro photography of a human face, 8k
resolution, razor-sharp focus. Enhance micro-details, crystal clear
optical sharpness. Authentic natural skin texture. The subject
smoothly turns their head towards the lens, ending the video with
a perfectly straight, direct frontal view portrait. Exact identity
preservation maintained throughout the temporal motion, maintaining
eye contact in the final frames.''
\end{quote}

\paragraph{Moderate blur ($\sigma \in \{8, 10\}$).}
\begin{quote}\small\itshape
``Cinematic ultra-high definition portrait of a human, 8k resolution,
raw photography. Exact structural reconstruction, authentic skin
texture with natural real-world anomalies. Professional studio
lighting. Dynamic temporal consistency: the camera naturally
transitions to lock onto a strict frontal view. In the final frames,
the face is perfectly aligned symmetrically towards the viewer,
ensuring highly detailed, unmodified facial structure from the front.''
\end{quote}

\paragraph{Severe blur ($\sigma \in \{12, 15\}$).}
\begin{quote}\small\itshape
``Hyper-realistic forensic-level restoration of a human face, 8k
resolution, extremely precise facial geometry. Strict structural
coherence, exact mapping of original facial proportions. Photorealistic
raw texture. The video dictates a deliberate motion ending in a static,
straight-on frontal view. The final frames strictly lock into a direct
frontal portrait, resolving all facial geometries symmetrically without
altering the original identity, neutral unedited appearance.''
\end{quote}

The prompt design follows two guiding principles: (i)~enforce
\emph{identity and structural preservation} by explicitly requesting
that the original facial geometry and biological attributes remain
unchanged, and (ii)~prescribe a \emph{motion trajectory} that
culminates in a frontal view, maximizing the likelihood of at least
one frame being suitable for face recognition.

The key hyperparameters used for all experiments are summarized in
Table~\ref{tab:i2v_params}.

\begin{table}[h]
\centering
\caption{TurboDiffusion I2V generation hyperparameters.}
\label{tab:i2v_params}
\begin{tabular}{l l}
\toprule
\textbf{Parameter} & \textbf{Value} \\
\midrule
Base model               & Wan2.2-A14B (14B parameters) \\
Resolution               & 720p (adaptive) \\
Number of frames ($T$)   & 81 \\
Inference steps           & 4 \\
Attention type            & SageSLA \\
SLA top-$k$ ratio         & 0.1 \\
Sampling method           & ODE \\
Weight quantization       & INT8 (quant\_linear) \\
Random seed               & 0 (fixed) \\
Noise-regime boundary     & 0.9 \\
\bottomrule
\end{tabular}
\end{table}

%------------------------------------------------------------
\subsubsection{Optimal frontal frame selection}
\label{subsubsec:frame_selection}

A critical component of the I2V pipeline is the extraction of
a single \emph{best frame} $I^{*}$ from the generated video
$\mathbf{V}$, representing the highest-fidelity frontal face.
We emphasize that this selection is performed \emph{without
access to the reference image}, making it applicable in real-world
scenarios where no clean reference is available at inference time.

\paragraph{InsightFace-based selection (full-video scan).}
For the InsightFace recognition system, which provides explicit
3D head-pose estimates via its Simple3D module, we scan
\emph{all $T$ frames} and select the most frontal frame using
a weighted pose-error heuristic:
\begin{equation}
    \varepsilon(I_t) = \alpha \cdot |\psi_t|
                     + \beta  \cdot |\theta_t|
                     + \gamma \cdot |\phi_t|,
    \label{eq:pose_error}
\end{equation}
where $(\psi_t, \theta_t, \phi_t)$ denote the yaw, pitch, and
roll angles of the detected face in frame $I_t$, and the weights
$(\alpha, \beta, \gamma) = (1.5,\; 1.0,\; 0.5)$ reflect the
relative importance of each axis---yaw is penalized most
heavily, as it has the greatest impact on facial feature visibility
for recognition.  The optimal frame is:
\begin{equation}
    I^{*} = \arg\min_{I_t \in \mathbf{V}} \varepsilon(I_t),
    \quad \text{subject to} \quad
    c(I_t) > \tau_{\text{det}},
    \label{eq:optimal_frame}
\end{equation}
where $c(I_t)$ is the detection confidence and
$\tau_{\text{det}} = 0.6$ is the minimum confidence threshold.
If multiple faces are detected in a frame, the largest face
(by bounding-box area) is selected.

\paragraph{DeepFace-based selection (landmark proxy).}
Since the DeepFace library does not provide explicit 3D pose
angles, we estimate frontality from 2D facial landmarks.
Denoting the left and right eye positions as
$\mathbf{p}_{\text{le}}, \mathbf{p}_{\text{re}}$ and the nose
tip as $\mathbf{p}_{\text{nose}}$, we define proxy angles:
\begin{align}
    \hat{\psi} &= \frac{|x_{\text{nose}} - \bar{x}_{\text{eye}}|}
                       {\|\mathbf{p}_{\text{re}} - \mathbf{p}_{\text{le}}\|_2}
                  \times 30,
    \label{eq:proxy_yaw} \\[3pt]
    \hat{\phi} &= \left|\arctan\!\left(
                  \frac{y_{\text{re}} - y_{\text{le}}}
                       {x_{\text{re}} - x_{\text{le}}}
                  \right)\right|,
    \label{eq:proxy_roll} \\[3pt]
    \hat{\theta} &= \left|
                    \frac{y_{\text{nose}} - \bar{y}_{\text{eye}}}
                         {\|\mathbf{p}_{\text{re}} - \mathbf{p}_{\text{le}}\|_2}
                    - 0.65
                    \right| \times 30,
    \label{eq:proxy_pitch}
\end{align}
where $\bar{x}_{\text{eye}} = (x_{\text{le}} + x_{\text{re}})/2$
is the horizontal eye midpoint, and $0.65$ is the expected
vertical nose-to-eye ratio for a frontal face.  The same weighted
combination as Eq.~\eqref{eq:pose_error} is then applied with
$(\hat{\psi}, \hat{\theta}, \hat{\phi})$ replacing
$(\psi, \theta, \phi)$.

\paragraph{Multi-criteria scoring (face extraction).}
For visual inspection and comparison-image generation, we also
employ a finer-grained scoring function that considers five
criteria simultaneously:
\begin{equation}
    S(I_t) = \underbrace{0.25 \cdot c_t}_{\text{detection}}
           + \underbrace{0.15 \cdot r_t}_{\text{face size}}
           + \underbrace{0.35 \cdot f_t}_{\text{frontality}}
           + \underbrace{0.15 \cdot \min\!\bigl(\tfrac{s_t}{500}, 1\bigr)}_{\text{sharpness}}
           + \underbrace{0.10 \cdot \tfrac{t}{T}}_{\text{position}},
    \label{eq:multi_score}
\end{equation}
where $c_t$ is detection confidence, $r_t$ is face-area ratio
(face area / frame area), $f_t = \max(0,\; 1 - |\hat{\psi}_t| / 45)$
is the frontality score, $s_t$ is the Laplacian variance
(sharpness), and $t/T$ encodes a slight preference for later
frames (which are more likely to show the frontal view given
our prompt design).

%------------------------------------------------------------
\subsubsection{Batch inference pipeline}
\label{subsubsec:batch_pipeline}

The full I2V restoration pipeline is orchestrated by a batch
inference script that automates the following workflow for each
experimental condition $(v, \sigma)$:

\begin{enumerate}
    \item \textbf{Input parsing.}\;
          Read the degraded-image CSV file for condition
          $(v, \sigma)$, listing all $N$ subjects with their
          degraded test image paths.

    \item \textbf{Prompt selection.}\;
          Select the severity-adaptive prompt $\mathcal{P}_\sigma$
          based on the blur level $\sigma$.

    \item \textbf{Video generation.}\;
          For each subject $i$, invoke
          \texttt{wan2.2\_i2v\_infer.py} with the degraded image
          $B_i(\sigma)$, prompt $\mathcal{P}_\sigma$, and the
          hyperparameters in Table~\ref{tab:i2v_params}, producing
          a video $\mathbf{V}_i$ saved as MP4.

    \item \textbf{Frame selection \& evaluation.}\;
          The evaluation scripts (Section~\ref{sec:evaluation})
          scan each generated video, apply the pose-based frame
          selection (Eq.~\ref{eq:optimal_frame}), extract the
          512-d embedding from the selected frame, and compute
          cosine similarity against the reference embedding.
\end{enumerate}

In total, the pipeline processes
$250 \times 6 \times 2 = 3{,}000$ video generation tasks across
all subjects, blur levels, and view types.  With the
\texttt{--skip\_existing} flag, the pipeline supports resumable
execution for overnight batch runs.

%------------------------------------------------------------
\subsubsection{Temporal redundancy hypothesis}
\label{subsubsec:temporal_redundancy}

The I2V approach offers a fundamental advantage over single-image
restoration: \emph{temporal redundancy}.  Among $T = 81$ generated
frames, the probability of at least one frame exhibiting
high-fidelity identity features is substantially higher than
from a single restoration attempt.  Formally, if we model the
probability that a single generated frame preserves identity above
a threshold as $p$, then the probability that at least one of $T$
independent frames exceeds the threshold is:
\begin{equation}
    P(\text{at least one success}) = 1 - (1-p)^T.
    \label{eq:temporal_redundancy}
\end{equation}

Even for $p = 0.05$ (only 5\% per-frame success rate), with
$T = 81$ frames the probability of at least one successful frame
is $1 - 0.95^{81} \approx 0.984$.\;  This multi-trial mechanism
is unavailable to single-image methods, which are constrained to
$T = 1$.

Furthermore, the video diffusion model implicitly learns the
manifold of plausible facial appearances under continuous viewpoint
changes.  Even when conditioned on a severely degraded side-profile
input, the model generates geometrically consistent frontal views
by leveraging the learned 3D facial prior, effectively performing
\emph{implicit 3D-aware face hallucination} through temporal
motion generation.
